\chapter{Búsquedas}

\section{Búsqueda binaria}
La función \texttt{lower\_bound} es una herramienta fundamental en la biblioteca estándar de C++ para realizar búsquedas eficientes en contenedores ordenados.

\subsection{Función \texttt{lower\_bound}}
\begin{itemize}
	\item \textbf{Complejidad}: $O(\log n)$ para contenedores de acceso aleatorio como \texttt{vector}
	\item \textbf{Comportamiento}: Devuelve un iterador al primer elemento que \textbf{no es menor} que el valor dado (primer elemento $\geq$ al valor buscado)
	\item \textbf{Sintaxis}: \texttt{lower\_bound(iterador\_inicio, iterador\_fin, valor)}
\end{itemize}

\section{Problema de ejemplo}

\subsection{Descripción}
El aula se representa como una línea unidimensional con celdas desde $1$ hasta $n$. Inicialmente hay $m$ profesores y David en posiciones distintas. En cada movimiento:
\begin{enumerate}
	\item David se mueve a una celda adyacente o se queda
	\item Todos los profesores se mueven simultáneamente a una celda adyacente o se quedan
\end{enumerate}

David es atrapado cuando algún profesor está en la misma celda. Todos actúan óptimamente: David maximiza el número de movimientos hasta ser atrapado, mientras los profesores lo minimizan.

\subsection{Solución}
Para cada consulta con la posición $a_i$ de David:
\begin{enumerate}
	\item Encontrar los dos profesores más cercanos (izquierda y derecha)
	\item Calcular la mitad de la distancia entre ellos
	\item David se moverá al punto medio para maximizar el tiempo
\end{enumerate}

\subsection{Código de Solución}
\begin{lstlisting}[style=cppstyle, caption={Solución del problema}]

	void solve() {
		ll n; cin >> n;
		ll m, q; cin >> m >> q;
		vi b(m);
		fore(i, 0, m) cin >> b[i];
		vi a(q);
		fore(i, 0, q) cin >> a[i];
		
		sort(all(b));
		
		fore(i, 0, q) {
			ll izq = 0, der = 0;
			ll res = 0;
			auto it = lower_bound(all(b), a[i]);
			
			if (it == b.begin()) {
				res = (*it) - 1;
			} else if (it == b.end()) {
				--it;
				res = n - (*it);
			} else {
				der = *it;
				--it;
				izq = *it;
				ll dist_der = der - a[i];
				ll dist_izq = a[i] - izq;
				res = (dist_der + dist_izq) / 2;
			}
			pri(res);
		}
	}
\end{lstlisting}

\section{Uso de \texttt{lower\_bound} para obtener índices}

\subsection{Obtención del índice}
Cuando se utiliza \texttt{lower\_bound} sobre un contenedor, podemos obtener la posición (índice) del elemento encontrado mediante aritmética de punteros:

\begin{lstlisting}[style=cppstyle, caption={Obtención del índice con lower\_bound}]
	void solve() {
		ll n; cin>>n; ll k; cin>>k;
		vi a(n); fore(i, 0, n) cin>>a[i];
		
		while(k--) {
			ll val; cin >> val;
			auto it = lower_bound(all(a), val);
			ll res = it - a.begin();  // Obtenemos el indice
			pri(res + 1);  // +1 para indice base 1
		}
	}
\end{lstlisting}

\section{Búsqueda binaria en la respuesta}
\subsection{Caso general}
\begin{lstlisting}[style=cppstyle, caption={Busqueda binaria en la respuesta}]
	bool can(int mid) {
		// Esta funcion debe responder si con mid como respuesta
		// es posible cumplir con la condicion del problema
		
		// Implementar la logica del chequeo
		
		return true;
	}
	
	int main(){
		// Busqueda binaria
		int low = 0, high = 1e9; // Ajustar limites segun el problema
		int ans = -1;
		
		while (low <= high) {
			int mid = (low + high) / 2;
			
			// si queremos el minimo que cumple
			if (check(mid)) {
				ans = mid;        
				high = mid - 1;
			} else {
				low = mid + 1;
			}
			
			// Si queres el maximo que cumple:
			if (check(mid)) {
				ans = mid;
				low = mid + 1;
			} else {
				high = mid - 1;
			}
		}
		
		cout << ans << endl;
	}
	
\end{lstlisting}

\subsection{Ejemplo 1: Problema de Sumas}
\begin{lstlisting}[style=cppstyle, caption={Problema de sumas con búsqueda binaria}]
	ll n;
	bool can(ll x, vll &a) {
		ll cont = 0; // Sup. disponibles
		fore(i, 0, n) {
			if (x < a[i]) return false; 
			cont += x - a[i];
		}
		return cont >= x;
	}
	
	void solve() {
		cin >> n; 
		vll a(n); 
		fore(i, 0, n) cin >> a[i]; 
		
		ll l = 0, r = 1e14, ans = 0;
		while (l <= r) {
			ll mid = (l + r) >> 1; 
			if (can(mid, a)) {
				ans = mid; 
				r = mid - 1;
			} else {
				l = mid + 1;
			}
		} 
		pri(ans);
	}
\end{lstlisting}

\subsection{Ejemplo 2: Problema de Porcentajes}
\begin{lstlisting}[style=cppstyle, caption={Problema de porcentajes con tickets}]
	ll k, x, y, a, b, n;
	vll pr;
	
	bool can(ll st) {
		vll p(st);
		for (int i = a - 1; i < st; i += a) p[i] += x;
		for (int i = b - 1; i < st; i += b) p[i] += y;
		sort(all(p), greater<ll>());
		ll cur = 0;
		fore(i, 0, st) cur += p[i] * pr[i] / 100;
		return cur >= k;
	}
	
	void solve() {
		pr.clear();
		cin >> n;
		fore(i, 0, n) {
			ll ai;
			cin >> ai;
			pr.pb(ai);
		}
		sort(all(pr), greater<int>()); 
		cin >> x >> a >> y >> b >> k;
		
		ll l = 0, h = n, ans = -1;
		while (l <= h) {
			ll mid = l + (h - l) / 2;
			if (can(mid)) {
				ans = mid;
				h = mid - 1;
			} else {
				l = mid + 1;
			}
		}
		pri(ans);
	}
\end{lstlisting}

\subsection{Ejemplo 3: Problema de Recursos}
\begin{lstlisting}[style=cppstyle, caption={Problema de gestión de recursos}]
	ll k, n; 
	vi a, b; 
	
	bool can(ll mid) {
		ll ka = k;
		fore(i, 0, n) {
			ll totact = a[i] * mid; 
			if (totact > b[i]) {
				ll rest = totact - b[i];
				if (rest <= ka) {
					ka -= rest;
				} else {
					return false;
				}
			}
		}
		return true;
	}
	
	void solve() {
		cin >> n >> k; 
		a.resize(n);
		b.resize(n);
		fore(i, 0, n) cin >> a[i]; 
		fore(i, 0, n) cin >> b[i]; 
		
		ll low = 0, high = 2e9 + 10;
		ll res = -1;
		
		while (low <= high) {
			ll mid = (low + high) / 2;
			if (can(mid)) {
				res = mid;
				low = mid + 1; // Buscamos el MAXIMO
			} else {
				high = mid - 1;
			}
		}
		pri(res);
	}
\end{lstlisting}

\subsection{Consideraciones Importantes}
\begin{itemize}
	\item \textbf{Límites}: Definir correctamente \texttt{l} y \texttt{r} según el problema
	\item \textbf{Función can}: Debe ser eficiente (generalmente $O(n)$ o $O(n \log n)$)
	\item \textbf{Overflow}: Usar \texttt{mid = l + (r - l) / 2} para evitar overflow
	\item \textbf{Tipo de búsqueda}: Decidir si se busca el mínimo o máximo que cumple
\end{itemize}

\subsection{Complejidad}
La complejidad total es $O(T \times \log(\text{rango}) \times \text{complejidad de can})$, donde:
\begin{itemize}
	\item $T$: número de casos de prueba
	\item \texttt{rango}: diferencia entre \texttt{r} y \texttt{l}
	\item \texttt{can}: complejidad de la función de verificación
\end{itemize}


